\chapter*{Résumé}
\addcontentsline{toc}{chapter}{Résumé}

Ce rapport présente la conception et la réalisation d'un jeu de \emph{bataille navale}
contre l'ordinateur, écrit en C11 avec la bibliothèque graphique GTK~3.

Le projet est né en 2023 d'un travail en binôme pour le cours \emph{Initiation aux interfaces
    graphiques}. La première version tenait dans un unique fichier de \num{2666} lignes : le jeu
fonctionnait, mais les règles, l'affichage et les réactions aux clics y étaient mêlés, et
chacun des trois niveaux de difficulté était une copie du même code.

La seconde version, objet de ce document, conserve le jeu et en refait les fondations. Elle
applique jusqu'au bout l'organisation \emph{modèle / vue / contrôleur} enseignée dans le
cours :

\begin{itemize}
    \item un \textbf{moteur de jeu} autonome, bibliothèque C11 sans aucune dépendance, sans
          allocation dynamique et sans état global ;
    \item \textbf{deux interfaces} pour ce même moteur, l'une graphique (GTK~3), l'autre en
          mode terminal --- preuve que le modèle ignore réellement sa vue ;
    \item \textbf{trois adversaires} dont la force est \emph{mesurée} sur des milliers de
          parties reproductibles : il faut en moyenne \num{96.2}, \num{65.8} et \num{55.3}
          tirs aux niveaux facile, moyen et difficile pour couler une flotte ;
    \item une \textbf{chaîne de qualité} complète : 54 cas de test unitaires, des tests de
          propriétés rejouant \num{2250} parties, un test de bout en bout de l'interface,
          l'analyse statique et dynamique du code, l'intégration et la livraison continues.
\end{itemize}

Le rapport décrit l'architecture retenue, les structures de données et les algorithmes ---
placement aléatoire d'une flotte, déductions tirées de la règle de non-contact, carte de
probabilité du niveau difficile ---, puis la démarche de vérification. Il se termine par un
bilan chiffré comparant les deux versions.

\vspace{1.2em}
\noindent\textbf{Mots-clés :} langage C, GTK, modèle-vue-contrôleur, programmation
modulaire, intelligence artificielle de jeu, tests, intégration continue.

\chapter*{Abstract}
\addcontentsline{toc}{chapter}{Abstract}

\begin{otherlanguage}{english}
    This report describes the design and implementation of a \emph{Battleship} game against the
    computer, written in C11 with the GTK~3 toolkit.

    The project started in 2023 as a pair assignment: a single \num{2666}-line file in which
    rules, display and event handling were intertwined. Version~2 keeps the game and rebuilds its
    foundations on a strict model / view / controller architecture: a dependency-free game engine,
    a graphical and a terminal front-end sharing it, three computer opponents whose strength is
    measured over thousands of reproducible games (\num{96.2}, \num{65.8} and \num{55.3} shots on
    average to sink a fleet), and a complete quality pipeline --- unit and property tests, an
    end-to-end interface test, static and dynamic analysis, continuous integration and delivery.
\end{otherlanguage}
